\section{An invariant graph and its exact finite gradient flow}
\label{sec:invariant-dynamics}

The stationary metric correction of Section~\ref{sec:projected-he}
does not define an invariant graph for metric evolution. We
construct such a graph directly from the heat estimates of
Section~\ref{sec:spatial-entry}. Fix a finite regularity order
$m\ge4$. The constants and the chosen graph may depend on $m$.

\begin{theorem}\label{thm:invariant-graph}
After a common shrink there is a chosen $C^m$ map
$\psi_{p,w}:\mathfrak h\to QX_p$, with $\psi_{p,w}(0)=0$,
such that
\begin{equation}\label{eq:graph-derivative-bounds}
 \sup_S\|D_S^j\psi_{p,w}(S)\|_X\le C_mr^2
                          \quad(0\le j\le m).
\end{equation}
For $j\le m-1$ the same bound holds for
$L_0D_S^j\psi$ in $Y_p$. The graph is invariant for
the extension defined below and locally invariant for actual
HYM. Its complete finite field $f_{p,w}$ satisfies
$\|f_{p,w}\|_{C^m}\le C_mr^2$.
On a fixed inner ball, a trajectory of this field reconstructs
the actual metric flow
\begin{equation}\label{eq:graph-reconstruction}
             h(t)=h_0\exp\{S(t)+\psi_{p,w}(S(t))\}.
\end{equation}
The graph points there are spatially smooth on each fixed fibre.
At zero the tangent map $J_0=I+D\psi_{p,w}(0)$ identifies
$\mathfrak h$ with the complete seven-dimensional small
spectral space of $L$, and
\begin{equation}\label{eq:graph-linear-intertwining}
                         J_0Df_{p,w}(0)=-LJ_0.
\end{equation}
All exact kernel directions are retained.
The construction is jointly $K$-equivariant in $w$ and $S$.
\end{theorem}

\subsection{A smooth extension of the first-jet equation}

Choose smooth scalar cutoffs $\theta,\chi$ taking values in
$[0,1]$, equal to one on $[0,1]$ and zero on $[4,\infty)$.
For a fixed small $a>0$, put
\[
 \widehat S=\chi(|S|^2/a^2)S,\qquad
 \vartheta(v)(x)=\theta\bigl((|v(x)|^2+|D_0v(x)|^2)/a^2\bigr).
\]
The extended nonlinearity is
\begin{equation}\label{eq:jet-extension}
 \widetilde N_w(S,v)(x)=\vartheta(v)(x)
                                  N_w(\widehat S+v)(x).
\end{equation}
Its right side is evaluated using the pointwise formula
\eqref{eq:Nw-definition}, with $D_0\widehat S=0$.
Choose $a$ so that the support of the cutoff lies strictly
inside that formula's pointwise domain, and extend by zero
outside. This uses a smooth cutoff in finite-dimensional
first-jet variables. It requires no smooth bump function
of the $C^1$ norm.

The bounded jet map $v\mapsto(v,D_0v)$ and the pointwise
chain rule make $\widetilde N_w:\mathfrak h\times QX_p\to Y_p$
globally smooth and bounded, with uniform fixed-order
derivative bounds. On the cutoff support the drift term
has size $O(ra)$ and the quadratic term has size $O(a^2+r^2a^2)$.
A cutoff derivative contributes at most $C/a$.
For $r\le a$ this proves
\begin{equation}\label{eq:extended-small-lipschitz}
 \|D\widetilde N_w\|\le\ell=C(a+r),\qquad
 \sup_S\|D_S^j\widetilde N_w(S,0)\|_Y\le C_jr^2.
\end{equation}
The second bound follows because every occurrence of a
parallel variable in \eqref{eq:Nw-definition} retains two
connection factors or the zeroth-order coefficient.
Also $\widetilde N_w(0,0)=0$.
Set
$\mathcal A(S,v)=\Pi\widetilde N_w(S,v)$ and
$\mathcal B(S,v)=Q\widetilde N_w(S,v)$, absorbing the
uniform projection norms into $\ell$. The extended system is
\begin{equation}\label{eq:extended-system}
 \dot S=\mathcal A(S,v),\qquad
 \dot v=-L_0v+\mathcal B(S,v).
\end{equation}
It agrees with actual HYM when $|S|<a$ and $\|v\|_X<a$.
The integrable smoothing kernel and global Lipschitz bounds
give unique global forward mild solutions on each fibre.
All cutoffs use invariant norms, so the extension is
jointly equivariant under parallel unitary conjugation.

\subsection{Backward histories and finite differentiability}

Let $k(t)=M(1+t^{-1/2})e^{-\sigma t}$ and
$K_\beta=\int_0^\infty k(t)e^{\beta t}\,\dd t$ for
$0\le\beta<\sigma$. For continuous histories on $(-\infty,0]$
use the norm
\[
 \|(S,v)\|_\beta=
       \sup_{t\le0}e^{\beta t}(|S(t)|+\|v(t)\|_X).
\]
Put $\eta=\sigma/(2(m+2))$. Reduce $a$ and then $r_0$
so that
\begin{equation}\label{eq:history-contractions}
 \ell(\beta^{-1}+K_\beta)<\tfrac12
 \quad\bigl(\eta\le\beta\le(m+1)\eta\bigr),
 \qquad \ell K_0<\tfrac12.
\end{equation}
For any endpoint $S_0$, solve
\begin{equation}\label{eq:backward-fixed-point}
 \begin{aligned}
 S(t)&=S_0-\int_t^0\mathcal A(S(s),v(s))\,\dd s,\\
 v(t)&=\int_{-\infty}^t e^{-(t-s)L_0}
                               \mathcal B(S(s),v(s))\,\dd s.
 \end{aligned}
\end{equation}
The two Lipschitz bounds are $\ell/\beta$ and
$\ell K_\beta$. Boundedness of the nonlinearities shows
that the map sends each indicated history space into itself:
the finite output grows at most linearly backwards, and
the normal output is bounded. Continuity of the latter
in $X_p$ follows by splitting off a short terminal interval
and using integrability of $k$. Thus there is a unique
fixed history, identical in all these weighted spaces.
Its normal part satisfies
\[
 \sup_{t\le0}\|v(t)\|_X
 \le K_0\{Cr^2+\ell\sup_{t\le0}\|v(t)\|_X\}\le Cr^2.
\]
Define $\psi(S_0)=v(0)$. The zero history gives $\psi(0)=0$.

We record the differentiability argument to account for
products of growing histories. The first variation solves
the linearized fixed-point equation with initial term
$(\delta S_0,0)$. Its integral operator has norm below
$1/2$ in each weight. Differences of fixed histories are
$O(|\delta S_0|e^{-\eta t})$, so the quadratic Taylor
remainder is $O(|\delta S_0|^2e^{-2\eta t})$.
Inverting in the $2\eta$ space proves actual
differentiability. At order $j$, the inhomogeneous terms
have the form
\[
 D^q(\mathcal A,\mathcal B)(S,v)
                 [Z_{i_1},\ldots,Z_{i_q}],
 \qquad q\ge2,\quad i_1+\cdots+i_q=j,
\]
with each $i_k<j$. Their weights add to $j\eta$.
The same inverse works there, and Taylor's formula with
the next derivative in the spare weight $(j+1)\eta$
proves differentiability and continuity. All these weights
are below $\sigma$. Evaluation at zero proves $C^m$
regularity of $\psi$ with uniform derivative bounds.

For the sharper $r^2$ bound, bounded mixed derivatives,
the normal bound, and \eqref{eq:extended-small-lipschitz}
give $D_S^j(\mathcal A,\mathcal B)=O(r^2)$ along every
fixed history. In the first normal variation, the finite
coefficient is therefore $O(r^2)$ and the normal coefficient
is small; absorption makes that variation $O(r^2)$.
At higher orders, a term with only finite variations has
an $O(r^2)$ coefficient, while every other term contains
an $O(r^2)$ normal variation. Induction in the weighted
equations proves \eqref{eq:graph-derivative-bounds}.
Uniqueness also proves joint equivariance.

\subsection{Invariance, domains and linearization}

Shift a fixed history by a finite time, attaching the forward
mild solution when necessary. Its negative tail still has
bounded normal component and at most linear finite growth.
The shifted normal integral has no term at minus infinity.
It therefore solves \eqref{eq:backward-fixed-point} with
its new endpoint. Uniqueness gives $v(t)=\psi(S(t))$.
Thus the graph is invariant and its field is
\begin{equation}\label{eq:finite-field}
                       f(S)=\mathcal A(S,\psi(S)).
\end{equation}
The chain rule and the preceding estimates give
$\|f\|_{C^m}\le C_mr^2$. This bounded field has a
complete finite flow in both time directions.

On the graph the mild equation and differentiation in
$X_p$ imply the generator-domain identity
\begin{equation}\label{eq:graph-domain-identity}
             L_0\psi(S)=\mathcal B(S,\psi(S))-D\psi(S)f(S).
\end{equation}
Indeed divide the mild formula at time $t$ minus its
initial value by $t$ and let $t\downarrow0$ in $Y_p$.
Strong continuity makes the integral quotient converge
to $\mathcal B$ and the left derivative equal $D\psi f$.
This is the definition of the generator domain.
Differentiating the right side and using closedness of
$L_0$ gives the claimed domain bounds through order $m-1$.

On a fixed smooth fibre, ordinary elliptic regularity now
gives $\psi,D\psi\in W^{2,q}$ for every finite $q$,
then $C^{1,\alpha}$. Equation~\eqref{eq:graph-domain-identity}
gives $C^{2,\alpha}$ regularity. In an inner finite ball,
a short graph trajectory stays where the cutoffs are
inactive, hence is a classical solution of the original
first-jet equation. Interior fixed-fibre parabolic
regularity makes it spatially smooth. These higher
regularity statements have no asserted uniform constants.
They justify the actual reconstruction
\eqref{eq:graph-reconstruction}.

At the HE origin $DN_w(0)=-(L-L_0)$.
Differentiating invariance gives
\eqref{eq:graph-linear-intertwining}.
The seven-dimensional image of $J_0$ is invariant under
the self-adjoint $L$ and lies in its operator domain.
Its restriction is self-adjoint for the induced mass;
the bound on $Df(0)$ puts all its eigenvalues in $[0,Cr^2]$.
The complementary gap in Theorem~\ref{thm:complete-bundle-spectrum}
then identifies this image with the whole small spectral
space. Since $\Pi J_0=I$, its coordinates are exactly
those of the inverse restricted projection.

If $S$ is in the constant actual parallel kernel, then
$DS=0$ and $h_0e^S$ is HE. Its constant history gives
$\psi(S)=f(S)=0$ in the inner ball. Conversely a finite
equilibrium reconstructs an HE metric. Integrating the
squared-distance identity in the proof of
\cite[Lemma 6.4.5]{Kobayashi} along the fibrewise geodesic
between two stationary HE metrics forces their logarithmic
difference to be $D$-parallel. Theorem~\ref{thm:complete-bundle-spectrum}
identifies that difference with a constant element of
$\mathfrak h$. Thus the local equilibrium set is exactly
the actual parallel kernel ball. This completes the proof
of Theorem~\ref{thm:invariant-graph}.

\subsection{The finite metric, energy and full Hessian}

Put $Z(S)=S+\psi(S)$ and $Z_Sa=a+D\psi(S)a$.
The pullback of the invariant $L^2$ metric on positive
Hermitian metrics is
\begin{equation}\label{eq:finite-flow-metric}
 g_{{\rm fin},S}(a,b)=\int_{X_p}\tr
             \{J_{Z(S)}(Z_Sa)J_{Z(S)}(Z_Sb)\}\,\dd V_{g_p}.
\end{equation}
Although the factors need not be $h_0$-Hermitian, they
are Hermitian for the current metric, so this is positive.
On the small ball it is uniformly equivalent to the
$L^2$ trace pairing of $Z_Sa,Z_Sb$, hence to $M$ since
$\Pi Z_Sa=a$.
Define
\begin{equation}\label{eq:finite-donaldson-energy}
 \mathcal E_{\rm fin}(S)=
 2\int_0^1\!\int_{X_p}
       \tr\{\mathrm i\Lambda F_{h_0e^{tZ(S)}}Z(S)\}
                                      \,\dd V_{g_p}\,\dd t.
\end{equation}
This is twice Donaldson's functional relative to $h_0$.
Its first variation, with $\dd V=\omega^3/3!$, is
\cite[Section 6.3, equations (6.3.25)--(6.3.27)]{Kobayashi}
\[
 \dd\mathcal D_h(\dot h)=\int_{X_p}
                      \tr(\mathrm i\Lambda F_h\,h^{-1}\dot h)\,\dd V.
\]
On the inner region where the graph carries the actual HYM
equation, graph invariance and its coefficient $-2$ give
\begin{equation}\label{eq:exact-finite-gradient}
                f=-\operatorname{grad}_{g_{\rm fin}}\mathcal E_{\rm fin}.
\end{equation}
This metric describes variations of metrics on one fixed
holomorphic bundle. The harmonic metric on varying bundle
moduli will be constructed separately.

\begin{corollary}\label{cor:relative-finite-hessian}
At an HE parameter $w=\chi_p(z)$, the exact finite mass
and Hessian satisfy
\begin{equation}\label{eq:finite-mass-hessian}
 \begin{split}
 g_{{\rm fin},0}(S,T)
   &=M(S,T)+\langle D\psi(0)S,D\psi(0)T\rangle_2,\\
 |g_{{\rm fin},0}-M|&\le C|z|^4M,\\
 (1-C|z|)E_{p,z}&\le\operatorname{Hess}_0\mathcal E_{\rm fin}
                                             \le(1+C|z|)E_{p,z}.
 \end{split}
\end{equation}
In particular the Hessian has the exact kernel of $E_{p,z}$,
and its bilinear error is bounded by
$C|z|\sqrt{E_{p,z}(S)E_{p,z}(T)}$.
\end{corollary}

\begin{proof}
At the critical origin the exact forms are
$N(S,T)=\langle J_0S,J_0T\rangle_2$ and
$H(S,T)=\langle LJ_0S,J_0T\rangle_2$.
Orthogonality gives the mass formula and its $O(r^4)$ bound.
The lower bound in \eqref{eq:whole-bundle-form} gives
$H\ge(1-Cr)E_{p,w}$.
Their kernels agree because $J_0$ is the inverse of
$\Pi$ on the full small spectral space and the actual
kernel consists of the constant stabilizer directions.
On the quotient by that common kernel, let $\bar N,\bar M$
be the quotient masses (the infima over representatives).
They are relatively $1+O(r^4)$ equivalent.
The positive generalized eigenvalues of $(H,N)$ are exactly
the positive small eigenvalues of $L$. Therefore
\eqref{eq:seven-eigenvalues} gives
\[
 \frac{\det\bar H}{\det\bar E_{p,w}}
  =\frac{\prod_{\lambda_j>0}\lambda_j}
         {\prod_{\mu_j>0}\mu_j}
                           \frac{\det\bar N}{\det\bar M}
  =1+O(r).
\]
Every eigenvalue of $\bar E_{p,w}^{-1/2}\bar H
\bar E_{p,w}^{-1/2}$ is at least $1-Cr$, and their
product is $1+O(r)$. Since the dimension is at most seven,
each is at most $1+Cr$. If the quotient has dimension
zero the assertion is immediate. This proves the full
relative form comparison without a lower bound for a
positive weak rate.
Finally equivariance of $\chi_p$ gives
$\delta_wS=d\chi_p(z)\delta_zS$, with derivative
$I+O(|z|)$ in the uniformly positive arrow norm.
Hence $E_{p,w}=(1+O(|z|))E_{p,z}$ as forms, including
kernels. Polarization on the quotient proves the
bilinear estimate.
\end{proof}

The exact linear generator is $-g_{{\rm fin},0}^{-1}H$.
The relative Hessian statement does not replace the
nonlinear field \eqref{eq:finite-field} by its quadratic
approximation on arbitrarily long time intervals.
